% captions_supplement.tex: generated by gmpe-bench; graphics paths are relative to results/paper/

\begin{figure}[htbp]
\centering
\adjustimage{max width=\textwidth}{supplement/figS1_random_leakage_inflation.pdf}
\caption{Leakage inflation: test RMSE under a record-level random split versus the event-grouped split (full profile), averaged over folds and IMs. Labels give the RMSE increase when events are held out, i.e. the optimism of random-split evaluation.}
\label{fig:figS1_random_leakage_inflation}
\end{figure}

\begin{figure}[htbp]
\centering
\adjustimage{max width=\textwidth}{supplement/figS2_period_coverage_heatmap.pdf}
\caption{Coverage at nominal 90\% per intensity measure and split scheme (full profile), averaged over models and folds: (a) split conformal, (b) HSCP. The colour scale is monotonic and colour-blind safe (cividis); values are printed in each cell. Split schemes: Evt = event-grouped 5-fold; Stn = station-grouped 5-fold; Reg = leave-one-region-out; Temp = temporal (train before 2008, test 2008 onward); $M_w{\geq}7$ = magnitude hold-out (test events with $M_w \geq 7$); CL1$\rightarrow$2 = mainshocks (CL1) for training, aftershocks (CL2) for testing.}
\label{fig:figS2_period_coverage_heatmap}
\end{figure}
